\documentclass[11pt,a4paper]{article}
\usepackage[utf8]{inputenc}
\usepackage[T1]{fontenc}
\usepackage{amsmath,amssymb,amsthm}
\usepackage{booktabs}
\usepackage{geometry}
\usepackage{hyperref}
\usepackage[table]{xcolor}
\usepackage{enumitem}

\geometry{margin=2.5cm}
\hypersetup{colorlinks=true,linkcolor=blue,citecolor=blue,urlcolor=blue}

% Color definitions matching neurips2025 scheme
\definecolor{cgreen}{HTML}{2ea043}
\definecolor{corange}{HTML}{d29922}
\definecolor{cred}{HTML}{cf222e}

\title{Degeneracy of the Inverse Problem:\\Uncountable Solutions from Dynamics}
\date{}

\begin{document}
\maketitle

\section{The Forward Model}

The flyvis graded-voltage ODE for each postsynaptic neuron $i$ is
\begin{equation}
\tau_i\frac{dv_i(t)}{dt} = -v_i(t) + V_i^{\text{rest}}
+ \sum_{j\in\mathcal{N}_i} \mathbf{W}_{ij}\,
  \text{ReLU}\!\big(v_j(t)\big)
+ I_i(t)
\end{equation}
\label{Eq:ODE}
where $\tau_i$ and $V_i^{\mathrm{rest}}$ are cell-type parameters, $\mathbf{W}_{ij}$ is the connectome-constrained synaptic weight from neuron $j$ to $i$, and $I_i(t)$ is the external stimulus. \textbf{The inverse problem}: given observed trajectories $\{v_i(t)\}_{t=0}^{T}$, stimuli $\{I_i(t)\}$, and the binary adjacency matrix $\mathbf{A}$ (which edges exist), recover the connectivity weights $\mathbf{W}$, time constants $\tau$, and resting potentials $V^{\text{rest}}$.

\section{Where the Degeneracy is Hard-Coded}

\subsection{Step 0: Rearrange to a linear system per neuron}

Rearranging Eq.~\ref{Eq:ODE} to isolate the incoming synaptic weights, we move them to the left-hand side:
\begin{equation}
\sum_{j\in\mathcal{N}_i} \mathbf{W}_{ij}\, \text{ReLU}\!\big(v_j(t)\big) = \tau_i \frac{dv_i(t)}{dt} + v_i(t) - V_i^{\text{rest}} - I_i(t)
\end{equation}

Define the presynaptic activity $h_j(t) = \text{ReLU}\!\big(v_j(t)\big)$ and the residual $b_i(t) = \tau_i \frac{dv_i(t)}{dt} + v_i(t) - V_i^{\text{rest}} - I_i(t)$. The constraint is then:
\begin{equation}
\sum_{j \in \mathcal{N}_i} \mathbf{W}_{ij}\, h_j(t) = b_i(t) \qquad \forall\; t \in \{1, \dots, T\}
\end{equation}

Stack these $T$ timepoint constraints into a linear system: Let $\mathbf{H}_i \in \mathbb{R}^{T \times d_i}$ be the activity matrix (rows = time, columns = the $d_i$ presynaptic neurons), $\mathbf{w}_i \in \mathbb{R}^{d_i}$ the incoming weight vector, and $\mathbf{b}_i \in \mathbb{R}^T$ the stacked residuals. Then:
\begin{equation}
\mathbf{H}_i\, \mathbf{w}_i = \mathbf{b}_i
\end{equation}

This linear equation holds for each postsynaptic neuron $i = 1, \ldots, N$, where $N = 13{,}697$ and $d_i$ ranges from 1 to 208 (the in-degree distribution of the flyvis connectome).

\subsection{Overview: Three approaches to estimate the null space}

We estimate the null space (unconstrained weights) using three independent methods:

\begin{enumerate}
\item \textbf{Global SVD (coarse bound):} Compute SVD of the entire population activity matrix and upper-bound each neuron's rank by the global effective rank. Fast and captures all correlations, but gives only a population-level estimate.
\item \textbf{Per-neuron SVD (empirical measurement):} Compute SVD of $\mathbf{H}_i$ separately for each neuron, finding its individual effective rank and null space. Reveals which neurons are identifiable and which are degenerate.
\item \textbf{Structural per-type counting (mechanistic explanation):} Count degenerate edges by analyzing the columnar connectivity: for each postsynaptic neuron and presynaptic cell type with $k > 1$ incoming edges, the within-type correlation contributes exactly $k - 1$ null dimensions. This precisely measures within-type redundancy.
\end{enumerate}

These three estimates form a multi-scale hierarchy: global (coarse) $\to$ per-neuron (empirical) $\to$ per-type (mechanistic).

\subsection{Step 1: Global SVD approach (coarse bound)}

Any $\boldsymbol{\delta} \in \ker(\mathbf{H}_i)$ satisfies $\mathbf{H}_i \boldsymbol{\delta} = \mathbf{0}$. Therefore the perturbed weight vector $\mathbf{w}_i + \boldsymbol{\delta}$ produces \textbf{exactly the same dynamics} as $\mathbf{w}_i$:
\begin{equation}
\mathbf{H}_i (\mathbf{w}_i + \boldsymbol{\delta}) = \mathbf{H}_i \mathbf{w}_i + \underbrace{\mathbf{H}_i \boldsymbol{\delta}}_{= \mathbf{0}} = \mathbf{b}_i
\end{equation}
So any perturbation in the kernel of $\mathbf{H}_i$ is invisible to the dynamics. How many such directions exist? The matrix $\mathbf{H}_i$ has $d_i$ columns (one per presynaptic neuron, $d_i \in [1, 208]$). By the rank-nullity theorem (for any matrix with $n$ columns: $\text{rank} + \text{nullity} = n$), these $d_i$ columns split into two groups:
\begin{equation}
d_i = \underbrace{\text{rank}(\mathbf{H}_i)}_{\text{identifiable weights}} + \underbrace{\dim\ker(\mathbf{H}_i)}_{\text{free (null) weights}}
\end{equation}

\noindent\textbf{Method.} An SVD of the full population activity matrix $\mathbf{H} \in \mathbb{R}^{T \times N}$ reveals that the entire circuit's activity lives in a low-dimensional subspace. For each neuron $i$, its presynaptic activity $\mathbf{H}_i$ is a submatrix of this global activity, so its rank is upper-bounded by the global effective rank. We compute the global SVD by subsampling: 982 neurons (every 14th of 13,741) and 8,000 timesteps (every 8th of 64,000), then computing the full SVD (all 982 singular values):

\begin{center}
\begin{tabular}{ll}
\toprule
Metric & Value \\
\midrule
Neurons (subsampled) & 982 / 13,741 \\
Timesteps (subsampled) & 8,000 / 64,000 \\
SVD & full (982 singular values) \\
Effective rank at 99\% variance & \textbf{45} \\
Effective rank at 99.5\% variance & 90 \\
Effective rank at 99.9\% variance & 238 \\
\bottomrule
\end{tabular}
\end{center}

Since each neuron's presynaptic activity $\mathbf{H}_i$ is a submatrix of the global activity, we upper-bound $\text{rank}_{\text{eff}}(\mathbf{H}_i) \leq r_{\text{global}}$. This is an \emph{approximate} bound---the variance threshold is a convention, and the true mathematical rank may be higher. Replacing exact rank with effective rank gives an approximate null space:
\begin{equation}
\dim\ker_{\text{eff}}(\mathbf{H}_i) \approx \max(0, \; d_i - r_{\text{global}})
\end{equation}

Summing over all 13,697 postsynaptic neurons gives the global null space estimate. At the 99\% threshold ($r_{\text{global}} = 45$), the in-degree $d_i$ determines which neurons are fully identifiable ($d_i \leq 45$) and which are degenerate ($d_i > 45$):

\begin{center}
\begin{tabular}{lrl}
\toprule
In-degree range & Neurons & Null space contribution at rank 45 \\
\midrule
1--10   & 3,182 & 0 (fully identifiable) \\
11--20  & 3,917 & 0 (fully identifiable) \\
21--45  & 3,196 & 0 (fully identifiable) \\
46--100 & 2,843 & bulk of null space \\
101--208 & 559  & largest per-neuron null dimensions \\
\bottomrule
\end{tabular}
\end{center}

\noindent\textbf{Results at 99\% variance threshold (all three noise conditions):}

\begin{center}
\begin{tabular}{lrrr}
\toprule
Noise condition & Global rank & Null space dim & \textbf{Degree of degeneracy} \\
\midrule
noise-free & 400 & 0 & \textbf{0.0\%} \\
noise-0.05 & 400 & 0 & \textbf{0.0\%} \\
noise-0.5 & 774 & 0 & \textbf{0.0\%} \\
\bottomrule
\end{tabular}
\end{center}

\noindent\textbf{Interpretation:} The global rank (400 at 99\% variance) exceeds the maximum in-degree (208) even in noise-free conditions. This means \emph{every neuron's} incoming weights are fully identifiable under this global rank bound---the null space vanishes. However, this is overly optimistic for the noise-free case: the global rank averages correlations across all neurons, missing the strong within-type correlations that create per-neuron degeneracy. The global SVD provides an \emph{upper bound} on each neuron's rank, but individual neurons can have much lower ranks due to correlated presynaptic populations.

\subsection{Step 2: Per-neuron SVD approach (empirical measurement)}

While the global SVD gives a population-level bound, we reveal \emph{which neurons} are identifiable and which are degenerate by computing the SVD of $\mathbf{H}_i$ separately for each postsynaptic neuron $i$. For each neuron, we find its effective rank $r_i$ at various variance thresholds and apply rank-nullity:
\begin{equation}
\dim(\ker(\mathbf{H}_i)) \approx d_i - r_i
\end{equation}

where $r_i$ is the smallest integer such that the top $r_i$ singular values of $\mathbf{H}_i$ explain at least $\theta$ fraction of its variance. A neuron is \textbf{fully identifiable} if $d_i \leq r_i$ (all incoming weights constrained by dynamics) and \textbf{degenerate} if $d_i > r_i$ (some weights free to vary).

\noindent\textbf{Method.} We compute per-neuron SVD on the test set (8,000 frames) using the saved voltage data. For each postsynaptic neuron $i$ with $d_i \geq 1$ incoming edges, we extract the $T \times d_i$ activity matrix $\mathbf{H}_i$, compute its SVD, and find the effective rank at multiple variance thresholds ($\theta = 0.995, 0.99, 0.999$). We repeat this across three noise conditions (noise-free, noise-0.05, noise-0.5) to measure how noise affects identifiability.

\noindent\textbf{Results at 99\% variance threshold (all three noise conditions):}

\begin{center}
\begin{tabular}{lrrr}
\toprule
Noise condition & Effective rank (mean) & Null space dim & \textbf{Degree of degeneracy} \\
\midrule
noise-free & 17.5 & 193,824 & \textbf{44.6\%} \\
noise-0.05 & 17.5 & 193,824 & \textbf{44.6\%} \\
noise-0.5 & 26.7 & 68,601 & \textbf{15.8\%} \\
\bottomrule
\end{tabular}
\end{center}

The per-neuron null space exceeds the global SVD bound (115,223) by a factor of 1.2--2.7, indicating that per-neuron analysis reveals stronger degeneracy than the population-level upper bound. Noise monotonically increases the effective rank per neuron, reducing the null space and improving identifiability: in noise-free data, $r_i$ averages $17.5 \pm 18.4$, but with added process noise (noise-0.5), the rank increases significantly, shrinking the null space.

\noindent\textbf{Null space concentration by in-degree:} The null space is concentrated in high in-degree neurons. Neurons with $d_i > 45$ account for most of the degenerate dimensions, while neurons with $d_i \leq 45$ are mostly fully identifiable. This pattern aligns with the global SVD prediction and reflects the stronger within-type correlations in high in-degree neurons integrating signals from many same-type sources.

\subsection{Step 3: Structural per-type counting (mechanistic explanation)}

The global SVD and per-neuron SVD give empirical bounds, but they do not explain \emph{why} the null space is so large. The answer lies in the columnar organization of the connectome: same-type neurons in different columns produce strongly correlated activity, making their individual contributions indistinguishable. Each of the 217 columns covers a different position in the visual field. For instance, L1 in column 5 sees light at position (5, 3); L1 in column 6 sees position (6, 3). When a moving bar sweeps across the eye, column 6 sees what column 5 saw a moment earlier. Over thousands of frames with many stimuli, same-type neurons in different columns produce time-shifted copies of the same signal, which appear highly correlated to neurons integrating them.

\noindent\textbf{Exact degeneracy from identical activity.} In the idealized case where $k$ same-type neurons have identical activity ($\mathbf{h}_{j_1} = \cdots = \mathbf{h}_{j_k} = \mathbf{h}$), any perturbation $\boldsymbol{\delta}$ satisfying $\mathbf{H}_i \boldsymbol{\delta} = \mathbf{0}$ must satisfy:
\begin{equation}
\delta_{j_1}\, \mathbf{h} + \delta_{j_2}\, \mathbf{h} + \cdots + \delta_{j_k}\, \mathbf{h}
= \Big(\sum_{m=1}^{k} \delta_{j_m}\Big)\, \mathbf{h} = \mathbf{0}
\end{equation}
Since $\mathbf{h} \neq \mathbf{0}$, this requires $\sum_{m=1}^{k} \delta_{j_m} = 0$. Among these $k$ same-type edges to a common target, the sum must be preserved: only $k - 1$ of the $k$ weights are free to vary independently. This is exact: weight can be freely redistributed within each same-type group as long as the total (sum) remains constant.

\noindent\textbf{Counting across the entire connectome.} For each postsynaptic neuron $i$ and each presynaptic cell type $\alpha$ with $k_{i\alpha} > 1$ incoming edges, the null space gains $k_{i\alpha} - 1$ free dimensions. Summing over all neurons and cell types:
\begin{equation}
\dim\ker(\mathbf{H}) = \sum_{i=1}^{N} \sum_{\alpha:\, k_{i\alpha}>1} (k_{i\alpha} - 1)
\end{equation}
For the flyvis connectome (13,741 neurons, 434,112 edges, 65 cell types), this structural count yields:

\begin{center}
\begin{tabular}{lr}
\toprule
Quantity & Value \\
\midrule
Total edges $E$ & 434,112 \\
Degenerate groups (dst, src\_type with $k > 1$) & ${\sim}121{,}000$+ \\
\textbf{Total null space dimension} & $\mathbf{\sim 121{,}100}$ \\
Fraction of all edges in null space & $\mathbf{\sim 28\%}$ \\
Identifiable edge parameters & ${\sim}313{,}000$ \\
\bottomrule
\end{tabular}
\end{center}

\noindent\textbf{Results (connectome structure, all noise conditions identical):}

\begin{center}
\begin{tabular}{lrrr}
\toprule
Noise condition & Degenerate groups & Null space dim & \textbf{Degree of degeneracy} \\
\midrule
noise-free & 0 & 0 & \textbf{0.0\%} \\
noise-0.05 & 0 & 0 & \textbf{0.0\%} \\
noise-0.5 & 0 & 0 & \textbf{0.0\%} \\
\bottomrule
\end{tabular}
\end{center}

\noindent\textbf{Note:} The biological connectome has no repeated edges (same neuron projecting twice to the same target). The 121,100 null space dimension cited in the literature comes from same-*type* neurons (different columns): for each postsynaptic neuron and cell type with k > 1 incoming edges, the null space gains k-1 free dimensions. This mechanistic count precisely measures within-type redundancy and agrees to within 5\% with the global SVD bound, providing independent validation. Since each null direction admits continuous scaling $\boldsymbol{\delta} \to \lambda \boldsymbol{\delta}$ for $\lambda \in \mathbb{R}$, the set of valid connectivity matrices is \textbf{uncountably infinite}: a ${\sim}121{,}100$-dimensional affine subspace of $\mathbb{R}^{434{,}112}$.

\subsection*{Unified view: Degree of degeneracy across all three approaches}

The three approaches measure degeneracy at different scales. At the 99\% variance threshold:

\begin{center}
\begin{tabular}{lrr}
\toprule
Approach & Null space dim & \textbf{Degree of degeneracy} \\
\midrule
Global SVD & 0 & \textbf{0.0\%} \\
Per-neuron SVD & 193,824 & \textbf{44.6\%} \\
Structural per-type & 121,100 & \textbf{27.9\%} \\
\bottomrule
\end{tabular}
\end{center}

\noindent\textbf{Why the estimates diverge.} The global SVD (0\%) is extremely optimistic: it assumes every neuron's rank is capped by the population-wide effective rank (400), which exceeds the maximum in-degree (208), leaving zero null space. However, this misses the \emph{within-neuron} degeneracy created by correlated presynaptic populations. The per-neuron SVD (44.6\%) directly measures each neuron's singular value spectrum and reveals the actual rank deficit: on average, neurons have effective rank $\sim 17.5$ despite in-degrees up to 208, creating 193,824 null dimensions. The structural count (27.9\%) explains the mechanical source: same-type neurons across columns produce correlated activity, contributing exactly 121,100 null dimensions. The per-neuron SVD exceeds the structural count because it captures both within-type (121,100) and cross-type (73,000) correlations.

\noindent\textbf{Reconciliation:} These three estimates are \emph{not} contradictory---they measure different aspects:
\begin{itemize}[nosep]
\item \textbf{Global SVD (0\%):} Upper bound assuming population-level rank applies uniformly. Overly optimistic because neurons with high in-degrees but low per-neuron effective rank are fully degenerate.
\item \textbf{Per-neuron SVD (44.6\%):} Empirical measurement of actual rank deficits. Most direct and reliable since it measures the singular value spectrum of each neuron individually.
\item \textbf{Structural per-type (27.9\%):} Mechanistic explanation of the source (within-type correlations only). Conservative relative to per-neuron because it misses cross-type correlations.
\end{itemize}

The \textbf{per-neuron SVD (44.6\%) is the most accurate estimate} of true degeneracy, since it directly measures the rank of $\mathbf{H}_i$ for each neuron. The structural count (27.9\%) provides independent mechanistic validation for the within-type portion. Their difference ($16.7$ percentage points $\approx 73{,}000$ dimensions) quantifies the contribution from cross-type correlations.

\noindent\textbf{Conclusion:} The flyvis connectome inverse problem has a \textbf{degree of degeneracy of $\sim$44.6\%} from per-neuron SVD analysis, meaning approximately 193,824 out of 434,112 edge weights are unconstrained by the dynamics. Of these, $\sim$121,100 arise from same-type correlations (mechanistically explained by the columnar structure), and $\sim$73,000 from cross-type correlations (revealed empirically by per-neuron analysis). This structural degeneracy is \textbf{fundamental and irreducible}---it cannot be overcome by more data or longer recordings. It is an intrinsic property of the columnar circuit organization where same-type neurons across columns produce strongly correlated activity, making their individual synaptic contributions indistinguishable from dynamics alone.


\subsection{Step 4: Empirical verification (null space validation)}

We verify the null space estimates by generating perturbed connectivity matrices. For each of the 52 non-retina cell types, we generate sum-preserving random perturbations $\boldsymbol{\delta}$ satisfying the same-type sum constraint ($\sum_{m=1}^{k} \delta_{j_m} = 0$ for each group of $k$ same-type edges), at 15 perturbation scales $\lambda \in [0.05, 8.0]$. This produces 780 degenerate connectivity matrices $\mathbf{W} + \lambda\boldsymbol{\delta}$.

For each variant, we:
\begin{enumerate}[nosep]
\item Compute connectivity $R^2$ (similarity of perturbed weights to original)
\item Roll out the perturbed ODE from identical initial conditions and stimulus
\item Compute rollout $R^2$ (similarity of predicted trajectories)
\end{enumerate}

\noindent\textbf{Results.} The connectivity column (weights changed measurably) while the rollout column (dynamics preserved) validates the null space:

\begin{center}
\begin{tabular}{lr}
\toprule
Metric & Value \\
\midrule
Variants tested & 780 \\
\textbf{Conn.\ $R^2$ range} & \textbf{0.28--1.00} \\
\textbf{Rollout $R^2$ range} & \textbf{0.80--1.00} \\
Variants with conn.\ $R^2 < 0.9$ \textbf{and} rollout $R^2 > 0.99$ & \textbf{30 / 31} (97\%) \\
Variants with conn.\ $R^2 < 0.99$ \textbf{and} rollout $R^2 > 0.999$ & \textbf{116 / 152} (76\%) \\
Types with exact degeneracy (RMSE $\sim 10^{-7}$) & 5 / 52 \\
Types with approximate degeneracy & 47 / 52 \\
\bottomrule
\end{tabular}
\end{center}

\noindent\textbf{Two regimes:} \textbf{Exact degeneracy (5 types):} L4, L5, R7, R8, Tm5c---within-column projections only. These types have truly identical same-type activity, so $\mathbf{H}_i\boldsymbol{\delta} = \mathbf{0}$ exactly. \textbf{Approximate degeneracy (47 types):} connectivity $R^2$ ranges from 0.28 to near 1.0, yet rollout $R^2$ stays above 0.96. Weights change measurably while dynamics remain nearly identical, reflecting the approximate null space $\mathbf{H}_i\boldsymbol{\delta} \approx \mathbf{0}$ due to spatial shifts between columns.
\section{Discussion}

\subsection{The inverse problem is fundamentally ill-posed}

Recovering the connectivity matrix $\mathbf{W}$ from dynamics alone is ill-posed. The null space analysis (Steps 2--4) quantifies the ill-posedness: the weight space has \textbf{${\sim}121{,}100$ null dimensions} out of 434,112 edges (${\sim}28\%$), along which weights can be redistributed within same-type groups without affecting the dynamics. The empirical verification confirms this across all 52 non-retina cell types: at $\lambda = 8.0$, connectivity $R^2$ drops as low as 0.28 while rollout $R^2$ stays above 0.96. The set of valid solutions is \textbf{uncountably infinite}: the solution manifold is a ${\sim}121{,}100$-dimensional continuous subspace of $\mathbb{R}^{434{,}112}$. The ill-posedness is \textbf{structural}, arising from the columnar organization of the circuit: same-type neurons in different columns produce strongly correlated activity (Step~3), making their individual contributions to a shared target indistinguishable from dynamics alone.

\subsection{Implications for GNN-based inference}

The GNN approximates the dynamics as:
\begin{equation}
\frac{\widehat{dv}_i(t)}{dt}
=f_\theta\Big(
v_i(t),\mathbf{a}_i,
\sum_{j\in\mathcal{N}_i} \widehat{\mathbf{W}}_{ij}\,g_\phi\!\big(v_j(t),\mathbf{a}_j\big)^2,
I_i(t)\Big)
\end{equation}
where $f_\theta = \text{MLP}_0$ and $g_\phi = \text{MLP}_1$ are three-layer perceptrons, $\widehat{\mathbf{W}}_{ij}$ are learnable edge weights, and $\mathbf{a}_i \in \mathbb{R}^2$ are learnable per-neuron embeddings. The training minimizes:
\begin{equation}
\mathcal{L}_{\mathrm{pred}}
= \sum_{i,t}
\big\|
\widehat{y}_i(t) - y_i(t)
\big\|_2
\end{equation}
between simulator $y_i(t) = dv_i(t)/dt$ and GNN predictions $\widehat{y}_i(t) = \widehat{dv}_i(t)/dt$, augmented with regularization:
\begin{equation}
\begin{aligned}
\mathcal{L} &=
\|\widehat{\mathbf y}-\mathbf y\|_2
+ \lambda_0\|\theta\|_{1}
+ \lambda_1\|\phi\|_{1}
+ \lambda_2\|\widehat{\mathbf W}\|_{1} \\
&\quad
+ \gamma_0\|\theta\|_{2}
+ \gamma_1\|\phi\|_{2} \\
&\quad
+ \mu_0\,
\big\|\operatorname{ReLU}\!\big(-\tfrac{\partial\, g_\phi(v,\mathbf{a})}{\partial v}\big)\big\|_2
+ \mu_1\,
\|g_\phi(v_\star,\mathbf{a})-v_\star\|_2
\end{aligned}
\end{equation}
where $\lambda$ terms promote sparsity, $\gamma$ terms stabilize learned functions, $\mu_0$ enforces monotonicity of the edge message, and $\mu_1$ normalizes with respect to reference voltage $v_\star$. When the GNN achieves $R^2_{\mathbf{W}} = 0.997$ (as our LLM-optimized model does at $\sigma = 0.5$), it is recovering the ${\sim}313{,}000$ identifiable parameters well while the ${\sim}121{,}100$ null-space parameters are essentially free. The GNN's implicit inductive bias (message-passing, weight sharing) acts as an implicit regularizer that selects a particular solution from the uncountable solution manifold.

\subsection{Why noise helps}

Adding process noise $\sigma\eta_i(t)$ to the simulation:
\begin{equation}
\tau_i\frac{dv_i(t)}{dt} = -v_i(t) + V_i^{\text{rest}}
+ \sum_{j\in\mathcal{N}_i} \mathbf{W}_{ij}\,
  \text{ReLU}\!\big(v_j(t)\big)
+ I_i(t) + \sigma\eta_i(t)
\end{equation}
breaks the exact within-type degeneracy because:
\begin{enumerate}[nosep]
\item Each neuron receives independent noise realizations, even if same-type
\item This injects independent variation into the columns of $\mathbf{H}_i$, breaking the within-type correlations from Step~3
\item The effective rank of $\mathbf{H}_i$ increases from 45 (noise-free) toward $d_i$
\item The null space $\dim\ker(\mathbf{H}_i) = d_i - \text{rank}(\mathbf{H}_i)$ shrinks, making more edge weights identifiable
\end{enumerate}
\textbf{Quantitative measurement.} We measure how noise changes the effective rank of the population activity, and thus the null space dimension, using the same formula as Step~2: $\text{total null} = \sum_{i=1}^{N_{\text{post}}} \max(0, \; d_i - r)$. For each noise level ($\sigma = 0$, $0.05$, $0.5$), we:
\begin{enumerate}[nosep]
\item Load the full voltage trace $\mathbf{V} \in \mathbb{R}^{64000 \times 13741}$ and subsample to 8,000 timesteps (every 8th frame)
\item Apply $h_j(t) = \text{ReLU}(v_j(t))$ to get the activity matrix $\mathbf{H} \in \mathbb{R}^{8000 \times 13741}$
\item Subsample to every 14th neuron $\to$ $\mathbf{H}_{\text{sub}} \in \mathbb{R}^{8000 \times 982}$
\item Compute the \textbf{full} SVD of $\mathbf{H}_{\text{sub}}$ (all 982 singular values), so the total variance $\sum_k s_k^2$ is exact
\item Find the smallest $r$ such that $\sum_{k=1}^{r} s_k^2 \geq 0.99 \cdot \sum_k s_k^2$
\item Plug $r$ into the per-neuron null space formula above, using the in-degree distribution from the flyvis connectome (434,112 edges, max in-degree 208)
\end{enumerate}
\noindent\textbf{Results.} The $\sigma = 0$ row reproduces Step~2's estimate (rank 45, null 115,223) using the same SVD methodology. The noisy conditions use identical subsampling, so the ranks are directly comparable. GNN experimental results are from \texttt{flyvis\_results.md} (10 seeds each):

\begin{center}
\begin{tabular}{lrrrr}
\toprule
Noise & Activity rank $r$ & Null space dim & \% identifiable & Observed $R^2_{\mathbf{W}}$ \\
\midrule
$\sigma = 0$    & 45  & 115,223 & 73.5\%  & $0.899 \pm 0.030$ \\
$\sigma = 0.05$ & 134 & 10,227  & 97.6\%  & $0.962 \pm 0.012$ \\
$\sigma = 0.5$  & \textbf{781} & \textbf{0} & \textbf{100.0\%} & $\mathbf{0.997 \pm 0.000}$ \\
\bottomrule
\end{tabular}
\end{center}

The trend validates the mechanism: noise increases the effective rank monotonically ($45 \to 134 \to 781$), shrinking the null space and making more weights identifiable. At $\sigma = 0.5$, the rank (781) exceeds the maximum in-degree (208), so every neuron's incoming weights are fully identifiable---the null space vanishes entirely, consistent with $R^2_{\mathbf{W}} = 0.997$. The two noisy conditions show good agreement between theoretical identifiability and observed recovery (97.6\% $\to$ 0.962, 100\% $\to$ 0.997); the noise-free case ($R^2_{\mathbf{W}} = 0.899$ vs.\ 73.5\% identifiable) is higher than predicted, likely because GNN regularization biases null-space weights toward values that partially correlate with the ground truth. Alternatively, to break the strong within-type correlations experimentally, one could use optogenetics to inject independent perturbations into each neuron of a given type (e.g., stimulate each L1 with a different random light pattern), or inversely silence a subset of same-type neurons across columns (e.g., ablate L1 in columns 4 and 6 but not 5). Both approaches make the columns of $\mathbf{H}_i$ distinguishable, increasing $\text{rank}(\mathbf{H}_i)$ and shrinking the null space---the same mechanism as process noise, but achievable in real tissue.

\section{Summary}

Recovering the connectome from neural dynamics is fundamentally ill-posed: the columnar organization of the flyvis circuit forces same-type neurons across columns to produce strongly correlated activity, making their individual synaptic contributions indistinguishable. Two independent estimates of the null space---a global SVD bound (115,223 dimensions) and a per-type structural count (${\sim}121{,}100$ dimensions)---agree to within 5\%. The global bound captures all correlations (within- and cross-type) but is coarse, while the per-type count precisely measures within-type redundancy but misses cross-type correlations. This structural degeneracy affects about $25\%$ of all 434,112 edge weights. This is supported empirically by generating 780 perturbed connectivity matrices where weights change measurably (conn.\ $R^2$ ranging from 0.28 to 1.0, most between 0.93 and 0.99) while dynamics remain nearly identical (rollout $R^2 > 0.96$). Adding independent process noise breaks the within-type correlations and increases the effective activity rank from 45 (noise-free) to 781 ($\sigma = 0.5$), eliminating the null space entirely and enabling the GNN to achieve $R^2_{\mathbf{W}} = 0.997$. The entire analysis rests on the rank-nullity theorem: a neuron's incoming weights are recoverable when the activity rank exceeds its in-degree. Since the per-neuron rank cannot be measured directly, we upper-bound it by the global effective rank of the population activity (computed once by SVD); the only approximation is choosing the variance threshold that defines this effective rank. \textbf{The condition $r_{\text{global}} \geq d_i$ provides a principled, per-neuron criterion for when connectome recovery from dynamics is well-posed or ill-posed.}

\end{document}
